\clearpage
\section*{Which new buildings do the permit rules miss, and which do they count twice?}
\textit{September 23, 2026. Agent-drafted data review requested by Jacob.}

The permit-based construction data (\texttt{tasks/prepare\_permit\_construction}) link each City of Chicago
new-construction permit issued 2006--2022 to the Assessor records of the building it authorized, and keep Assessor
buildings that no permit reaches as a separate route. Two groups looked like coverage failures: 2,043 Assessor
buildings from 2008 on with no permit, and 2,018 residential permits that reached no building. Most of both were
problems in how the data were read, not missing records. After general rule changes the file has 18,864 rows
and 14,033 density-eligible buildings (14,047 before). Of those, 11,562 are permitted (11,378 before) and 4,959 are
multifamily (4,951 before). The total barely moves because the gains among permits offset duplicates removed from
the Assessor-only route.

Three reading errors hid permits. The permit classifier treated notes about later permits
(\texttt{[SEE PERMIT \#\ldots\ TO CONVERT TO 4 DWELLING UNITS]}) and misspellings (\texttt{DWELING},
\texttt{6 UNITBUILDING}) as conversions or non-residential work. A garage rule dropped every permit that mentioned a
garage without a dwelling word, which removed 282 new buildings such as \texttt{ERECT 2 STORY MASONRY BLDG PER PLANS
WITH DETACHED GARAGE}. Accessory structures are now identified by the object of the first clause, and permits
for a new building of unstated use are kept when their own parcels receive a new residential record that no
residential permit reaches (114 measured, 105 eligible). A list of non-residential uses keeps hotel foundations,
office build-outs, dormitories and equipment buildings out. On the Assessor side, hotel rooms and care beds in the
commercial valuations were counted as apartments; those uses are now excluded.

Two Assessor practices hid buildings from the permits that authorized them. For 247 permits the Assessor recorded
the new building under the old year built (a three-flat still dated 1890). The parcel's floor area jumps after the
permit, so a parcel without a new record whose floor area rises by at least 25 percent within the construction lag,
and stays there, is measured in that year. Of these, 228 are eligible; most were first assessed one or two years
after the permit, and the units agree with the permit in all but 17. Second, a condominium declaration or subdivision
retires a parcel number, and the building reappears on new parcels. The permit was then measured on the old lot,
often at zero floor area mid-construction, while the finished building appeared again as a separate no-permit
building. A parcel-succession rule links a record on retired parcels to the new records on the same tax block within
150 feet, first assessed within a year of the retirement, when their units add up to its own and they succeed no
other record. It gives 105 permits their final measurements and removes about 500 Assessor-only rows that counted a
building twice, 209 of them in the eligible 2006--2007 group. Historical parcel centroids from every assessment year
1999--2025 (\texttt{tasks/download\_parcel\_centroids}) locate 18,851 of the 18,864 rows (18,163 before), which
the succession rule needs. The 98 commercial records that were never located are suburban and were dropped.

What remains open: 236 permits are followed by a measured permit on the same parcel or address and were not built
under the earlier permit (\texttt{later\_permit\_built}); 1,603 permits still reach no building. Of those, 272 are
cancelled, expired or revoked, 113 sit on parcels whose floor area changed too little for the rule, and 120 are in
the review queues. Samples of the rest include tax-exempt affordable and senior housing without Assessor
characteristics, pre-2021 commercial apartment buildings, and genuinely unbuilt permits. Another 1,763 Assessor
buildings from 2008 on have no permit (600 with three or more units); sampled permits at their parcels are mostly
conversions and gut rehabs, which the Assessor re-dates but which are not new construction. A six-year construction
window instead of four added only 52 eligible buildings and made some large buildings' measurements worse, so the
window stays at four years. No regressions have been run on these data.
